\documentclass[11pt]{article}
% Adaptado de C:/GitHub/Mestrado/Template/Template.tex.
\usepackage[T1]{fontenc}
\usepackage[utf8]{inputenc}
\usepackage[brazilian]{babel}
\usepackage{graphicx,booktabs,array,amsmath,amssymb,microtype}
\usepackage{geometry}
\usepackage{pdflscape,caption}
% Apresentação de artigo para leitura em tela, conforme preferência do autor.
\geometry{a4paper,margin=2.5cm,headheight=12pt,headsep=0.5cm}
\usepackage{setspace,parskip,csquotes}
\setstretch{1.05}
% Fonte padrão original e parágrafos sem recuo, separados pelo espaço de parskip.
\setlength{\parindent}{0pt}
\usepackage{titlesec}
\titleformat{\section}[hang]{\normalfont\large\bfseries}{\thesection}{\fontdimen2\font}{}
\titleformat{name=\section,numberless}[block]{\normalfont\large\bfseries\filcenter}{}{0pt}{}
\titleformat{\subsection}[hang]{\normalfont\normalsize\bfseries}{\thesubsection}{\fontdimen2\font}{}
\titlespacing{\section}{0pt}{0pt}{8pt}
\titlespacing{\subsection}{0pt}{10pt}{5pt}
\titlespacing{\paragraph}{0pt}{7pt}{0.7em}
\AtBeginDocument{%
  \setlength{\abovedisplayskip}{8pt plus 2pt minus 2pt}%
  \setlength{\belowdisplayskip}{8pt plus 2pt minus 2pt}%
  \setlength{\abovedisplayshortskip}{4pt plus 2pt}%
  \setlength{\belowdisplayshortskip}{6pt plus 2pt minus 2pt}%
}
\usepackage{fancyhdr}
\pagestyle{fancy}
\fancyhf{}
\fancyhead[R]{\footnotesize\thepage}
\renewcommand{\headrulewidth}{0pt}
\renewcommand{\footnoterule}{\kern-3pt\hrule width 5cm\kern2.6pt}
\makeatletter
\long\def\@makefntext#1{\parindent 0pt\noindent\hangindent=1.8em\hangafter=1\hb@xt@1.8em{\hss\@makefnmark}#1}
\makeatother
% Notas e legendas em tamanho 9; espaço simples nos objetos.
\DeclareCaptionFont{simples}{\setstretch{1}}
\captionsetup{font={footnotesize,simples},labelsep=endash,justification=raggedright,singlelinecheck=false,position=top,skip=6pt}
% Notas de figuras e tabelas, editáveis no LaTeX e com apresentação uniforme.
\newenvironment{notas}{%
  \par\smallskip\noindent\begin{minipage}{\linewidth}%
  \footnotesize\setstretch{1}\raggedright
  \setlength{\parindent}{0pt}\setlength{\parskip}{3pt}%
}{\end{minipage}}

\renewenvironment{abstract}{\par\medskip\noindent\textbf{\abstractname}\par\smallskip\small\setstretch{1}\noindent}{\par\medskip}
\interfootnotelinepenalty=10000
\widowpenalty=10000
\clubpenalty=10000
\displaywidowpenalty=10000
\usepackage[backend=biber,style=abnt]{biblatex}
\addbibresource{references.bib}
\setlength{\bibhang}{0pt}
\renewcommand*{\bibfont}{\normalfont\small\setstretch{1}}
\AtBeginBibliography{\setlength{\bibitemsep}{5pt}}
\DeclareFieldFormat{url}{\small\bibstring{urlfrom}\addcolon\addspace\url{#1}}
\DeclareFieldFormat{urldate}{\small\bibstring{urlseen}\addcolon\addspace #1}
\usepackage{hyperref}
\hypersetup{hidelinks}
\newcommand{\instituicao}{UNIVERSIDADE FEDERAL DE SANTA CATARINA}
\newcommand{\departamento}{Programa de Pós-Graduação em Economia}
\newcommand{\curso}{Mestrado}
\newcommand{\disciplina}{Desigualdade, Diversidade e Redes}
\newcommand{\titulo}{Rede de emissões de CO\textsubscript{2} no Brasil}
\newcommand{\autor}{Gabriel Pincelli Cintra}
\newcommand{\local}{Florianópolis}
\newcommand{\dataentrega}{18 de outubro de 2026}
\newcommand{\imprimirtitulodois}{
\begin{center}
  {\small\MakeUppercase{\instituicao}\par}
  {\small \departamento\par}
  {\small \disciplina\ (\curso)\par}
  \vspace{1em}{\LARGE\bfseries \titulo\par}
  \vspace{0.7em}{\normalsize \autor\par}
  \vspace{0.3em}{\small \local, \dataentrega\par}
\end{center}
\vspace{0.4em}}
\begin{document}
\imprimirtitulodois
\thispagestyle{empty}
\begin{abstract}
Este trabalho examina a distribuição das emissões de CO\textsubscript{2} entre atividades emissoras e setores de produtos finais no Brasil, complementando a comparação das contas pela produção e pela demanda. Combina a Matriz de Insumo-Produto brasileira de 2015, com 67 atividades, e coeficientes setoriais de emissão em uma rede dirigida e ponderada. Os volumes e a concentração caracterizam as relações intersetoriais, e o filtro de disparidade seleciona ligações pela importância relativa para a origem ou o destino. Os setores de energia e de água, esgoto e resíduos distribuem volumes semelhantes entre outros produtos finais, mas com concentrações distintas. Na construção, 39,5\% das emissões atribuídas aos produtos finais do setor são geradas em outras atividades, incluindo as produtoras de materiais e de energia. Uma mesma relação pode representar participações distintas para a origem e o destino, e a seleção retém proporções desiguais dos volumes setoriais. Esses resultados oferecem elementos para coordenar a redução das emissões nos setores fornecedores com a eficiência material e a aquisição de materiais de menor emissão na construção. A interpretação ambiental depende da validação dos coeficientes utilizados.
\end{abstract}
\noindent\textbf{Palavras-chave:} atribuição de emissões; insumo-produto; redes dirigidas e ponderadas; mitigação.

\clearpage
\input{secoes/introducao.tex}
% Sequência inspirada em Huang et al. (2019).
\clearpage
\section{Metodologia}
\input{secoes/metodologia_dados.tex}
\input{secoes/metodologia_modelo.tex}
\input{secoes/metodologia_rede.tex}
\clearpage
\input{secoes/resultados.tex}
\clearpage
\section{Discussão e implicações para a coordenação da mitigação}

No cenário analisado, energia e água/resíduos distribuem volumes intersetoriais semelhantes, com maior concentração dos destinos em água/resíduos, especialmente na administração pública. Na construção, a conta pela demanda reúne emissões próprias e de diversas origens. A decomposição das contas identifica as atividades emissoras e os produtos finais associados aos saldos setoriais, complementando as comparações entre setores, regiões e cadeias produtivas da literatura brasileira \parencite{imori2015,sanguinetAzzoni2024,montoyaEtAl2026}.

A leitura dirigida de \textcite{huang2019} permite examinar a participação de cada atividade nas atribuições de suas contrapartes. Construção representa 32,57\% das saídas externas de produtos de metal, enquanto essa origem responde por 4,42\% das entradas externas de construção. As duas perspectivas identificam o peso do destino para a origem e a contribuição da origem ao destino. O filtro de \textcite{serrano2009} explicita essa importância relativa sob seu modelo nulo, e as coberturas indicam o alcance em volume do conjunto selecionado.

Na construção, 39,5\% das emissões atribuídas aos produtos finais vêm de outras atividades e 60,5\% são geradas na própria atividade. Entre as origens externas estão produtos de minerais não metálicos, energia e produtos de metal. Essa composição relaciona decisões dos produtores, dos responsáveis pelos projetos e dos contratantes. A coordenação pode combinar a redução das emissões por unidade produzida com projetos que utilizem menos materiais para prestar o mesmo serviço, conforme as estratégias de eficiência material e os critérios de emissões incorporadas discutidos por \textcite[seções~5.1 e~6.4]{rissman2020}.

Para energia, mantidas a estrutura produtiva e a demanda do modelo, uma redução da intensidade emissora diminuiria as emissões atribuídas a diversos produtos finais. Esse alcance intersetorial relaciona a oferta energética de menor emissão às mudanças nos processos industriais discutidas por \textcite[seções~4.2 e~7.2]{rissman2020}.

A divulgação verificável das emissões dos materiais e os critérios de menor emissão nas compras públicas podem apoiar decisões de projeto e de aquisição na construção \parencite[seções~6.6 e~6.7]{rissman2020}. A aplicação desses instrumentos à construção requer informação ambiental dos produtos e fornecedores, avaliação de ciclo de vida e conhecimento da participação governamental na demanda do setor. A escolha dos instrumentos reúne essa informação às condições de aquisição, aos custos e aos critérios de distribuição dos esforços de mitigação.

A interpretação ambiental depende da compatibilidade dos coeficientes interpolados com a MIP de 2015. A validação da extensão ambiental exige conciliar a cobertura das emissões incluídas, a alocação nas 67 atividades e a base de preços da produção usada no denominador dos coeficientes. A revisão dos coeficientes pode alterar volumes, participações e relações selecionadas. A consistência contábil verifica a preservação das contas, enquanto a validação ambiental examina as emissões físicas representadas.

Como desdobramento da análise nacional, uma aplicação a Santa Catarina em 2018 poderia identificar os produtos finais associados às emissões de cada atividade e as origens que compõem cada conta pela demanda. Dados harmonizados e o recálculo das contas e dos perfis permitiriam examinar as contrapartes e a concentração das atribuições na fronteira estadual.

\section{Conclusões}

No cenário analisado, as emissões de energia e água/resíduos atribuídas a outros produtos finais superam as emissões provenientes de outras origens e atribuídas aos produtos finais de energia e água/resíduos. Em construção ocorre o inverso, ampliando sua conta pela demanda em relação à produção. A identificação das contrapartes esclarece esses saldos e distingue perfis de volumes próximos com concentrações diferentes, apesar da conectividade semelhante. Construção também mostra a necessidade de separar as emissões atribuídas aos produtos finais da própria atividade das relações intersetoriais. A leitura conjunta das contas e dos perfis permite compreender a composição das diferenças entre produção e demanda.

O backbone distingue relações destacadas para a origem, para o destino ou para ambos. Os casos mostram que essa importância pode ser assimétrica e que a estabilidade dos papéis examinados convive com mudanças na seleção e na cobertura.

Na cadeia da construção, os resultados oferecem elementos para considerar conjuntamente as emissões nos processos fornecedores, a eficiência no uso de materiais e critérios de menor emissão nos projetos e nas aquisições. A interpretação sustenta a discussão dessas estratégias, sem estabelecer a eficácia ou o custo de instrumentos. A aplicação a Santa Catarina requer contas e perfis próprios, e conclusões empíricas definitivas sobre o Brasil dependem da validação ou substituição da extensão ambiental.

\clearpage
\begingroup
\setstretch{1}
\raggedright
\printbibliography[title={Referências}]
\endgroup
\input{secoes/apendices.tex}
\end{document}
